%======================================================================
\chapter{Hagopian's decomposition}\label{ch:decomp}
%======================================================================

This chapter proves the following theorem, which is the first half of the proof that $M$ is hereditarily
indecomposable (the second half, Chapter~\ref{ch:noarc}, shows that the quotient it produces cannot exist).

\begin{theorem}[Hagopian 1984, Theorem~1, with $N=M$]\label{thm:h84}
Let $M$ be a homogeneous indecomposable plane continuum that is not hereditarily indecomposable. There is a
continuous surjection $q\colon M\to Y$ onto a compact metric space $Y$, with connected fibres none of which is $M$,
such that $Y$ is homogeneous and contains an arc.
\end{theorem}

\formal{\lean{HPC.Hagopian.Hagopian84Decomp} (the statement), \lean{HPC.Hagopian.hagopian84_decomp} (the proof).}

The route is Hagopian's 1984 paper \emph{Atriodic homogeneous continua}, which simplifies the decomposition of his
1976 paper and works for atriodic homogeneous continua in general. Since the argument is long and consists mostly of
careful bookkeeping with chains of open sets, this chapter gives the definitions and statements in full and the proofs
in outline; the complete arguments are in the Lean files \lean{HPC/Hagopian/H84*.lean} and
\lean{HPC/Hagopian/Hagopian84*.lean}. Throughout, $M$ is as in the standing hypotheses of Chapter~\ref{ch:homog},
and ``minimal admissible'' and ``layer'' are as in Chapter~\ref{ch:typelambda}.

\section{Rigidity of layers}

The following facts express that type $\lambda$ continua in $M$ fit together like arcs in a one-dimensional
manifold: where two of them overlap essentially, their layers agree.

\begin{lemma}\label{lem:rigid}
Let $(G,g)$ be minimal admissible in $M$.
\begin{enumerate}[label=(\alph*)]
\item \emph{(Staying inside.)} A subcontinuum $C$ of $M$ that meets $G$ and misses the end layers of $G$ lies in $G$.
\item \emph{(No arm at an interior layer.)} If $0<u<1$ and $A$ is a continuum with $\emptyset\ne A\cap G\subseteq\lay gu$, then
$A\subseteq\lay gu$.
\end{enumerate}
\end{lemma}

\begin{proof}
(a) is Step~1 of the proof of Theorem~\ref{thm:layer-invariant}, with $C$ in place of $h(Z)$. (b) Otherwise $\lay
g{[0,u]}$, $\lay g{[u,1]}$ and $A\cup\lay gu$ are three continua through the core $\lay gu$ that pairwise meet only in it,
and none lies in it; Corollary~\ref{cor:nested-or-arc}(b).
\end{proof}

\formal{\lean{HPC.Hagopian.subset_of_meets_avoids}, \lean{HPC.Hagopian.no_arm_at_interior_layer}.}

\begin{theorem}[Continuity of the decomposition; Hagopian 1976, Lemma~5]\label{thm:h76-lemma5}
If $(E,k)$ is minimal admissible in $M$, the decomposition of $E$ into layers is continuous: the map $s\mapsto\lay ks$ is
continuous for the Hausdorff metric.
\end{theorem}

\begin{proof}[Proof in outline]
Upper semicontinuity holds for every admissible map. Lower semicontinuity at an interior layer comes from
Theorem~\ref{thm:h76-lemma3}: if layers $\lay k{s_n}$ accumulated only on part of $\lay ks$, small homeomorphisms
taking points of $\lay ks$ to nearby points of $\lay k{s_n}$ (Proposition~\ref{prop:effros-nbhd}) would map $\lay ks$
onto $\lay k{s_n}$. At an end layer, Hagopian passes to a larger type $\lambda$ continuum in which it becomes an
interior layer: an \emph{extension} of $E$ away from the other end set (Hagopian 1976, Lemma~4: the extensions of
$E$ away from an end set are linearly ordered by inclusion and have no largest element; the formal proof supplies the
argument for ``no largest element'', which Hagopian leaves unargued). The layers of $E$ are layers of the extension.
\end{proof}

\formal{\lean{HPC.Hagopian.IsExtensionAway}, \lean{HPC.Hagopian.hagopian76_lemma4_chain},
\lean{HPC.Hagopian.hagopian76_lemma4}, \lean{HPC.Hagopian.hagopian76_lemma5}.}

\begin{corollary}[Restriction]\label{cor:restrict}
If $(L,d)$ is minimal admissible and $a<b$, then $\lay d{[a,b]}$ is of type $\lambda$, with a minimal admissible map
whose layers are the layers $\lay d{a+t(b-a)}$; its end sets are $\lay da$ and $\lay db$.
\end{corollary}

\formal{\lean{HPC.Hagopian.minimalAdmissible_restrict}.}

A subcontinuum $H$ of a type $\lambda$ continuum $E$ is \emph{essential} in $E$ if it meets more than one layer of $E$.

\begin{theorem}[Hagopian 1984, Lemma~6]\label{thm:h84-lemma6}
If $H$ is an essential subcontinuum of $(E,k)$, then $H=\lay k{[a,b]}$ for some $a<b$, $H$ is of type $\lambda$, and every
layer of the decomposition of $H$ is a layer of $k$.
\end{theorem}

\begin{proof}
By the interval lemma, $H\supseteq\lay k{(a,b)}$ where $[a,b]=k(H)$. By continuity of the decomposition
(Theorem~\ref{thm:h76-lemma5}) the layers over $a$ and $b$ are limits of layers over $(a,b)$, so $H=\lay k{[a,b]}$;
Corollary~\ref{cor:restrict} gives the minimal admissible map, and uniqueness (Theorem~\ref{thm:unique-decomp}) the last
statement.
\end{proof}

\formal{\lean{HPC.Hagopian.IsEssential84}, \lean{HPC.Hagopian.hagopian84_lemma6}.}

This replaces Thomas's Theorem~8, which Hagopian cites.

\begin{theorem}[Rigidity; Hagopian 1984, Lemma~7]\label{thm:layers-agree}
Let $(E,k)$ and $(G,d)$ be minimal admissible in $M$, and suppose $E\cap G$ meets two layers of $E$. Then either $E$ lies in a
single layer of $G$, or $E\cap G=\lay k{[a,b]}=\lay d{[a',b']}$ with $a<b$ and $a'<b'$, every layer of $k$ meeting $G$ is a layer
of $d$, and every layer of $d$ meeting $E$ is a layer of $k$. In particular (Hagopian's Lemma~7), if neither of $E,G$
contains the other and $G$ meets a layer $\lay kr$, then $\lay kr$ is a layer of $d$.
\end{theorem}

\begin{proof}[Proof in outline]
$K=E\cap G$ is a continuum (Corollary~\ref{cor:nested-or-arc}(a)). If $K$ lies in one layer of $G$ and $E\not\subseteq G$, then
$E$ is an arm at that layer and Lemma~\ref{lem:rigid}(b) applies (end layers are handled by passing to extensions,
Lemma~4 of 1976, in which they become interior). If $K$ meets two layers of each, Theorem~\ref{thm:h84-lemma6} applied in
$E$ and in $G$ shows that $K$ is a union of layers of both, and that the two decompositions of $K$ coincide by
uniqueness. The remaining case ($K$ meets two layers of $E$ and lies in one layer $Y$ of $G$) is excluded by passing to a
two-sided extension of $G$ in which $Y$ is interior, and applying Lemma~\ref{lem:rigid}(b) again. This route differs from
Hagopian's printed proof.
\end{proof}

\formal{\lean{HPC.Hagopian.layers_agree}, \lean{HPC.Hagopian.hagopian84_lemma7_aux}.}

\begin{remark}
Hagopian 1976, Lemma~6, and 1984, Lemma~11 (with $N=M$): the closure of the union of the extensions of $E$ away from an end
set is indecomposable (\formal{\lean{HPC.Hagopian.hagopian76_lemma6_X}, \lean{HPC.Hagopian.hagopian84_lemma11}.}) In
the 1976 paper this union plays the role that the decomposition $\mathfrak D$ below plays in 1984.
\end{remark}

\section{Chains through which a continuum runs straight}

We work with chains of five open subsets of $M$, indexed $L_0,\dots,L_4$ (Hagopian's $L_1,\dots,L_5$).

\begin{definition}\label{def:chain5}
A \emph{5-chain} is a sequence $L=(L_0,\dots,L_4)$ of open subsets of $M$ with $L_i\cap L_j\ne\emptyset$ iff $|i-j|\le1$. Its
\emph{middle} is $L'=L_1\cup L_2\cup L_3$. It is
\begin{itemize}
\item \emph{normal} if $\cl{L_i}\cap\cl{L_j}=\emptyset$ for $|i-j|>1$;
\item \emph{free} if $\fr\bigcup_iL_i\subseteq\fr(L_0\cup L_4)\setminus\cl{L'}$ (frontiers and closures in $M$).
\end{itemize}
A minimal admissible $(E,k)$ \emph{runs straight through} $L$ if (1) $E\subseteq\bigcup_iL_i$; (2) $\lay k0\subseteq L_0\setminus L_1$
and $\lay k1\subseteq L_4\setminus L_3$; (3) a layer meeting $\fr L_i$ lies in $\fr L_i$; (4) if $r<t$ and $\lay kr\cup\lay kt\subseteq
L_i$, then $\lay k{[r,t]}\subseteq L_i$. The chain is \emph{regular} if for each component $K$ of $L'$, $\cl K$ (a
\emph{strand}) is of type $\lambda$ and runs straight through $L$. An \emph{ordered refinement} of $L$ is a 5-chain $P$ with
$P_i\subseteq L_i$.
\end{definition}

\formal{\lean{HPC.Hagopian.IsChain5}, \lean{HPC.Hagopian.midUnion}, \lean{HPC.Hagopian.IsNormal5},
\lean{HPC.Hagopian.IsFree5}, \lean{HPC.Hagopian.RunsStraight}, \lean{HPC.Hagopian.IsRegular5},
\lean{HPC.Hagopian.IsOrderedRefinement5}.}

The chains are built from Effros neighbourhoods of pieces of $E$: for $S\subseteq M$ and $\delta>0$ let $\Eff(S,\delta)$ be the
set of points $h(x)$, $x\in S$, $h$ a $\delta$-homeomorphism; it is open by Proposition~\ref{prop:effros-nbhd}.

\begin{lemma}[Folds]\label{lem:fold}
Let $V_0,\dots,V_4$ be open subsets of $[0,1]$ covering it, with $V_i\cap V_j=\emptyset$ for $|i-j|>1$, $0\in V_0\setminus V_1$ and
$1\in V_4\setminus V_3$. If some $V_i$ is not an interval, there are $r<s<t$ and adjacent indices $p,q$ with $s\in V_p$ and
$r,t\in V_q\setminus V_p$ (a \emph{fold}).
\end{lemma}

\begin{proof}
Sweep from left to right. Suppose there is no fold. If $a\in V_p$, no point left of $a$ lies in a later set and no point $\ge
a$ in an earlier one, then $V_p$ is an interval, and the first point after $a$ outside $V_p$ lies in $V_{p+1}$, with the
same two properties. Starting with $a=0$, $p=0$, every $V_p$ is an interval. (If $V_p$ had a gap, the last point of the
complement before a point of $V_p$ and the first after it would lie in neighbours of $V_p$, and the two properties force both
into the same neighbour: a fold.)
\end{proof}

\formal{\lean{HPC.Hagopian.exists_fold}.}

\begin{theorem}[Hagopian 1984, Lemmas~12 and~13]\label{thm:h84-lemma13}
Let $(E,k)$ be minimal admissible in $M$, $\varepsilon>0$, $A$ a closed set missing $E$, and $B$ an open set containing $\lay k0\cup\lay
k1$. Then $E$ runs straight through a free normal regular 5-chain $P$ in $M\setminus A$ with $\cl{P_0\cup P_4}\subseteq B$, such that
for each component $K$ of the middle $P'$ there is an $\varepsilon$-homeomorphism $h$ for which $\cl K$ is an essential
subcontinuum of $h(E)$.
\end{theorem}

\begin{proof}[Proof in outline]
\emph{Lemma~13, the chain.} Choose $0=x_0<x_1<\dots<x_{15}=1$ with $\lay k{[0,x_3]}$ and $\lay k{[x_{12},1]}$ inside $B$, and let
$L_i=\Eff(\lay k{[x_{3i},x_{3i+3}]},\delta)$. For small $\delta$ these form a normal chain in $M\setminus A$ (pieces of non-adjacent
links are disjoint compact sets), and every $\delta$-translate $h(E)$ satisfies (1) and (2). Condition (3) uses the Effros
theorem with the room $\delta-\max_vd(v,h(v))$; condition (4) uses Lemma~\ref{lem:fold}: a non-interval among the parameter
sets $\set{s: h(\lay ks)\subseteq L_i}$ gives a fold, and folds are excluded by atriodicity
(Lemma~\ref{lem:rigid}(a)) and the gap bounds $\dist(\lay k{[0,x_j]},\lay k{[x_{j+1},1]})$. Two translates that meet have
their layers matched by Theorem~\ref{thm:h76-lemma3}. Regularity: a component $K$ of the middle contains a point $h(v)$,
$v\in E$; $\cl K\subseteq h(E)$ by Lemma~\ref{lem:rigid}(a), and it is essential in $h(E)$; by
Theorem~\ref{thm:h84-lemma6} it is of type $\lambda$ and runs straight.

\emph{Lemma~12, freeness.} Put $A'=L_0\setminus\cl{L_1}$ and $B'=L_4\setminus\cl{L_3}$. Moore's theorem on separating closed sets
(\vthree{Theorem~1.9}) applied in $M\setminus(A'\cup B')$ splits off the points that cannot be joined to $E$, and removing
them from the middle links gives an ordered refinement $P$ that is free, normal and regular, through which $E$ still runs
straight, and whose middle components are components of $L'$. The hypothesis of Moore's theorem comes from atriodicity
directly.
\end{proof}

\formal{\lean{HPC.Hagopian.WsetOf}, \lean{HPC.Hagopian.chainLinks84}, \lean{HPC.Hagopian.chainLinks84_chain},
\lean{HPC.Hagopian.chainLinks84_runsStraight}, \lean{HPC.Hagopian.chainLinks84_regular},
\lean{HPC.Hagopian.hagopian84_lemma12}, \lean{HPC.Hagopian.hagopian84_lemma13}.}

\section{The layers of a chain}

For a regular 5-chain $P$ let $\Omega(P)$ be the set of all layers of all strands of $P$, and $\Delta(P)$ the set of those
that miss $\cl{P_0\cup P_4}$.

\formal{\lean{HPC.Hagopian.Omega5}, \lean{HPC.Hagopian.Delta5}.}

\begin{proposition}\label{prop:delta}
Let $P$ be a regular 5-chain.
\begin{enumerate}[label=(\alph*)]
\item $\Delta(P)$ is a partition of $\bigcup_iP_i\setminus\cl{P_0\cup P_4}$.
\item If $P$ is free, there is $\eta>0$ with $d(z,y)\ge\eta$ for $z\in\bigcup\Delta(P)$ and $y\notin P'$.
\item \emph{(Lemma~14.)} If $P$ is free, $\Delta(P)$ is continuous: members through nearby points are Hausdorff close.
\item \emph{(Lemma~15.)} If $P$ is free and normal, there is $\varepsilon>0$ such that every $\varepsilon$-homeomorphism maps each member of
$\Delta(P)$ onto a member of $\Omega(P)$ inside $P'$.
\item \emph{(The arc.)} If $P$ is normal, every strand has a closed interval of layers $\lay c{[a,b]}$, $a<b$, all in $\Delta(P)$.
\end{enumerate}
\end{proposition}

\begin{proof}[Proof in outline]
(a) A layer through a point outside $\cl{P_i}$ misses $\cl{P_i}$ by (3); two members that meet come from the same strand,
where the decomposition is unique. (b) $\cl{\bigcup\Delta(P)}\subseteq P'$ by freeness, then compactness. (c) For $\varepsilon'=\min(\varepsilon,\eta)$,
an $\varepsilon'$-homeomorphism carrying $z$ to $z'$ maps the member through $z$ onto the member through $z'$ by
Theorem~\ref{thm:h76-lemma3} (its hypotheses are checked with $\eta$). (d) The image $h(Z)$ lies in one component $K'$ of
the middle; $h(\cl K)\cap\cl{K'}$ meets two layers of $h(\cl K)$ by Lemma~\ref{lem:rigid}(b), and
Theorem~\ref{thm:layers-agree} makes $h(Z)$ a layer of $\cl{K'}$. (e) The strand meets $P_2$, and the layer through such a
point lies in $P_2$, hence misses $\cl{P_0\cup P_4}$ by normality; the parameters of layers inside $P_2$ form an open set.
\end{proof}

\formal{\lean{HPC.Hagopian.delta5_partition}, \lean{HPC.Hagopian.delta5_dist_pos},
\lean{HPC.Hagopian.hagopian84_lemma14}, \lean{HPC.Hagopian.hagopian84_lemma15},
\lean{HPC.Hagopian.exists_interval_in_delta}.}

\begin{theorem}[Lemma~16$'$]\label{thm:lemma16p}
Let $P$ be a free normal regular 5-chain and $(F,d)$ minimal admissible with $\lay d0\cap P'=\emptyset$. Every layer of $F$ that meets
$P'$ is a layer of the strand of $P$ through any of its points in $P'$.
\end{theorem}

\begin{proof}[Proof in outline]
By Theorem~\ref{thm:layers-agree}, either the layer of $F$ is a layer of the strand, or the strand lies in a single layer of
$F$. The second alternative is impossible: following the layers of $F$ from the end layer $\lay d0$ (outside $P'$) to the
strand, a limit point argument in $\cl{P_2}$ produces an arm at an interior layer.
\end{proof}

\formal{\lean{HPC.Hagopian.strand_layer_mem_delta}, \lean{HPC.Hagopian.not_strand_subset_extension_layer},
\lean{HPC.Hagopian.strand_layer_of_extension}; Hagopian's Lemma~16 as printed is
\lean{HPC.Hagopian.hagopian84_lemma16}.}

\section{The decomposition}

Fix a type $\lambda$ continuum in $M$ (Theorem~\ref{thm:exists-lambda}) and a free normal regular chain $P$ through which it runs
straight (Theorem~\ref{thm:h84-lemma13} with $A=\emptyset$, $B=M$). The members of $\Delta(P)$ cover $\bigcup_iP_i$ except near the
end links; the rest of $M$ is covered by further chains.

\begin{proposition}[Bridges and the chains $P(Z)$]\label{prop:bridges}
Let $Z\in\Omega(P)\setminus\Delta(P)$. Then $Z$ lies in $\cl{P_0}$ or in $\cl{P_4}$ and misses $\cl{P_2}$. There is a type $\lambda$
continuum $J(Z)$ (a \emph{bridge}) containing $Z$ as a layer, missing $\cl{P_2}$, with both end layers in $\Delta(P)$. By
Theorem~\ref{thm:h84-lemma13} with $A=\cl{P_2}$ and $B=\bigcup\Delta(P)$, $J(Z)$ runs straight through a free normal regular chain
$P(Z)$ in $M\setminus\cl{P_2}$ with $\cl{P_0(Z)\cup P_4(Z)}\subseteq\bigcup\Delta(P)$, and then $Z\in\Delta(P(Z))$.
\end{proposition}

\formal{\lean{HPC.Hagopian.omega_not_delta_side}, \lean{HPC.Hagopian.exists_bridge84},
\lean{HPC.Hagopian.mem_delta_of_bridge_chain}, \lean{HPC.Hagopian.exists_chainZ84}.}

\begin{proposition}[The finite cover]\label{prop:finite-cover}
There are finitely many $Z_1,\dots,Z_n\in\Omega(P)\setminus\Delta(P)$ with
$M=\bigcup\Delta(P)\cup\bigcup_{i=1}^n\bigcup\Delta(P(Z_i))$.
\end{proposition}

\begin{proof}
$C=\cl{\bigcup\Delta(P)}\setminus\bigcup\Delta(P)$ is compact. Each point of $C$ lies in $P'\cap\cl{P_0\cup P_4}$, so the layer $Z$
through it of its strand is in $\Omega(P)\setminus\Delta(P)$, and the point lies in the open set $\bigcup\Delta(P(Z))$; finitely many of
these cover $C$. The union $U$ of $\bigcup\Delta(P)$ with them is open; it is closed because
$\cl{\bigcup\Delta(P(Z))}\subseteq\bigcup\Delta(P(Z))\cup\cl{P_0(Z)\cup P_4(Z)}$ (freeness) and $\cl{P_0(Z)\cup P_4(Z)}\subseteq\bigcup\Delta(P)$. As
$M$ is connected, $U=M$.
\end{proof}

\formal{\lean{HPC.Hagopian.exists_finite_cover84}, \lean{HPC.Hagopian.hagopian84_cover},
\lean{HPC.Hagopian.CoverData84}, \lean{HPC.Hagopian.exists_coverData84}.}

Hagopian argues instead that the union of the whole (uncountable) family is closed; the formal proof uses the finite form
above, and deduces the printed statement from it (\formal{\lean{HPC.Hagopian.sUnion_cover84_eq_univ}.})

Put $P(Z_0)=P$ and let $\mathfrak D=\bigcup_{i=0}^n\Delta(P(Z_i))$.

\begin{theorem}\label{thm:frakD}
\leavevmode
\begin{enumerate}[label=(\alph*)]
\item \emph{(Transfer.)} If $Z\in\Omega(P)$ contains a point $z$ of the middle of some $Q=P(Z_i)$, then $Z$ is the layer through $z$ of
the strand of $Q$ through $z$; if moreover $z\notin\cl{Q_0\cup Q_4}$, then $Z\in\Delta(Q)$.
\item $\mathfrak D$ is a partition of $M$ into proper subcontinua.
\item \emph{(Claims (5) and (6).)} $\Omega(P)\subseteq\mathfrak D$, and there is $\delta>0$ such that every $\delta$-homeomorphism maps members of
$\mathfrak D$ onto members of $\mathfrak D$.
\item $\mathfrak D$ is upper semicontinuous, and some compact $C\subseteq M$ has a continuous map onto $[0,1]$ whose fibres are the traces of
members of $\mathfrak D$ (an arc of members).
\end{enumerate}
\end{theorem}

\begin{proof}[Proof in outline]
(a) $Z=c^{-1}(r)$ for the strand map $c$; the strand meets $P_2$ in a layer $c^{-1}(s)$, which misses $\bigcup Q\ni z$ (as $Q$ lies in
$M\setminus\cl{P_2}$), so $s\ne r$; the piece of the strand between these two layers is of type $\lambda$
(Corollary~\ref{cor:restrict}) with one end layer missing $Q'$ and the other $Z$; apply Theorem~\ref{thm:lemma16p} to $Q$.
(b) Every point is covered (Proposition~\ref{prop:finite-cover}); two members that meet are equal, by (a) and
Proposition~\ref{prop:delta}(a), and in the last case by Theorem~\ref{thm:h76-lemma3}. (c) A member of $\Omega(P)\setminus\Delta(P)$ lies in
$\cl{P_0}$ or $\cl{P_4}$, so a point of it lies in some $\bigcup\Delta(P(Z_i))$; then (a). For (6), take $\delta$ below the constants of
Proposition~\ref{prop:delta}(b),(d) for each of the finitely many chains, and use (a) again. (d) Upper semicontinuity from
Proposition~\ref{prop:delta}(c); the arc from Proposition~\ref{prop:delta}(e).
\end{proof}

\formal{\lean{HPC.Hagopian.omega_layer_of_mem_mid}, \lean{HPC.Hagopian.omega_mem_delta_of_mem_mid},
\lean{HPC.Hagopian.CoverData84.members}, \lean{HPC.Hagopian.CoverData84.members_disjoint},
\lean{HPC.Hagopian.CoverData84.omega_subset_members}, \lean{HPC.Hagopian.CoverData84.invariant},
\lean{HPC.Hagopian.CoverData84.isJonesDecomp}.}

\section{The quotient}

Call an equivalence relation on $M$ a \emph{Jones decomposition} if it is upper semicontinuous, its classes are connected and different
from $M$, there is $\delta>0$ such that every $\delta$-homeomorphism respects it, and it has an arc of classes as in
Theorem~\ref{thm:frakD}(d). By Theorem~\ref{thm:frakD}, $\mathfrak D$ is one.

\formal{\lean{HPC.Hagopian.IsJonesDecomp84}.}

\begin{proof}[Proof of Theorem~\ref{thm:h84}]
Let $q\colon M\to Y=M/\mathfrak D$ be the quotient map. $Y$ is compact and metrizable (\vthree{Theorem~6.5}), $q$ is continuous and onto
with connected fibres, none equal to $M$.

\emph{$Y$ is homogeneous.} Let $\mathcal G$ be the set of finite compositions of $\delta$-homeomorphisms. Each $g\in\mathcal G$ respects $\mathfrak
D$ and so induces a homeomorphism of $Y$ (\vthree{Proposition~6.8}). The orbit of a point $x$ under $\mathcal G$ is open (it contains
$\Eff(y,\delta)$ around each of its points $y$, Proposition~\ref{prop:effros-nbhd}), and the orbits partition $M$, so each orbit is also
closed. As $M$ is connected, $\mathcal G$ acts transitively, and so does the group of induced homeomorphisms of $Y$.

\emph{$Y$ contains an arc.} The map $C\to[0,1]$ of Theorem~\ref{thm:frakD}(d) factors through $q(C)$ as a continuous bijection of a
compact space onto $[0,1]$.
\end{proof}

\formal{\lean{HPC.Hagopian.exists_comp_small_homeo}, \lean{HPC.Hagopian.hagopian84Decomp_of_package},
\lean{HPC.Hagopian.hagopian84_decomp}.}

\begin{remark}
In Hagopian's text the chain of points used for homogeneity, the positivity of several constants, the closedness of the union
$\bigcup\mathfrak B$, the choice of $\varepsilon$ in Lemma~15 and the existence of the arc in $M/\mathfrak D$ are stated without argument;
the formal proof supplies them as indicated. Several printed lemmas (7, 12, 15, 16) are proved by routes different from the
printed ones, and Lemma~16 is used in the stronger form Lemma~16$'$.
\end{remark}

\section*{Exercises}

\begin{exercise}
Prove Lemma~\ref{lem:fold} in detail, and show by an example that the conclusion can fail if $V_0$ and $V_2$ are allowed to meet.
\end{exercise}

\begin{exercise}
Let $Y$ be a compact metric space and $\mathcal G$ a set of homeomorphisms of $Y$ closed under composition such that the $\mathcal
G$-orbit of every point contains a neighbourhood of the point. Show that if $Y$ is connected, $\mathcal G$ acts transitively.
\end{exercise}
